\section{Problem Description}

I build a navigation framework that moves a mobile robot from a start location to a
goal location in a partially known 2D environment while avoiding obstacles
and respecting the robot's motion constraints. The robot is a differential-drive
platform: it cannot translate sideways and its linear and angular velocities are
bounded and acceleration-limited, so a valid plan must be more than a sequence of
free cells, it must be drivable.

The framework combines two complementary layers. A global planner searches the
whole environment for a complete, collision-free route, providing long-term guidance.
A local planner/optimiser then tracks that route while continuously scoring
candidate motions against the robot's dynamics and nearby obstacles.

\paragraph{Why global planning is necessary, and why local planning alone is not.}
A local planner only looks a short distance ahead, so it always tries to move straight
toward the goal. This works in open space but gets stuck when the goal is hidden behind an
obstacle. The clearest example is a dead end: if the robot drives into a U-shaped
wall, the goal is on the far side of it, so every ``move closer'' step just pushes the
robot into the wall. Escaping would mean first moving \emph{away} from the goal, which a
local planner never chooses --- so it stays trapped. A global planner avoids this by
looking at the whole map at once and routing around such traps before the robot even
starts. The two layers are complementary: the global planner decides which way to
go, the local planner how to drive there smoothly.
